% Authoritative LaTeX manuscript; Markdown is a generated compatibility projection.

Start with a small set of engineering models, not an organization-wide commitment to autonomy. Choose a part of the system whose structure, behavior, or obligations you understand well enough to describe and check. Identify what the models must represent, which questions they should answer, and how their claims can be validated. The first objective is modest: make enough engineering knowledge explicit to support a bounded assignment whose result can be evaluated.

Structural models are often a good starting point. Engineers already reason about components, dependencies, interfaces, and ownership boundaries, even when these relationships are not formally documented. Much of this structure can be recovered from source code, configuration, and existing documentation. Static analysis can help check whether the implementation agrees with the model. Begin with a few important relationships rather than attempting to describe the entire system. Distinguish the structure that exists from the structure you intend: recovering a dependency does not establish that it belongs there.

Before building machinery, examine how the selected models could improve engineering work. What must someone currently reconstruct from source, documents, or conversations? Which decisions recur? What must remain true while the implementation changes? Preserve the requirements and decisions that matter to acceptance, including qualitative obligations that still require judgment. Also identify choices that can safely remain open. A model that specifies every implementation decision can consume the very engineering effort delegation was meant to save.

Inventory the knowledge and enforcement you already have. A schema may reject invalid data, and a permission may prevent a dangerous action. A test may expose a regression, while a review may settle an ambiguous case. Determine which of these actually affects whether work proceeds. An adequate existing mechanism needs a place in the workflow, not a replacement bearing a new name. Where existing enforcement already provides adequate protection at a justified cost, additional system modeling may contribute little. Where people repeatedly recover the same architectural relation, a small representation may help before any new gate is needed.

\subsection{Choose a Modeling Approach}\label{heading:choose-a-modeling-approach}

Before accumulating models, establish how they will be interpreted and maintained. \autoref{2.2-constructing-and-using-engineering-models:heading:semantic-commitment-and-analysis} distinguishes semantic commitment from the means of storing and interacting with a model. Decide which questions and analyses the representation must support, then choose how engineers and agents will edit, query, visualize, and check it. \autoref{practicing-modeling-approaches} offers starting points for those decisions.

{\footnotesize
\begin{longtable}[]{@{}p{.25\linewidth}p{.31\linewidth}p{.34\linewidth}@{}}
\caption[Choosing modeling approaches and tools]{\emph{Modeling semantics and tooling serve different purposes.} These approaches can be combined; the rows are not a ranking of precision or value. Choose the combination whose analyses and uses justify its construction and upkeep.}\label{practicing-modeling-approaches}\\
\toprule
\textbf{Approach and examples} & \textbf{Semantic capability} & \textbf{Tooling role and practical use} \\
\midrule
\endhead
Informal models: sketches, ADRs, Markdown & Meaning is stated in prose and conventions, with interpretation left to readers. & Explore alternatives and record qualitative decisions before committing to a machine-interpretable form. \\
UML; PlantUML or Mermaid views~\citep{omg2017uml,plantumlDocumentation,mermaidDocumentation} & UML defines standard modeling constructs; a diagram renderer does not establish the semantics of every depicted relation. & Render architecture, behavior, or interaction views; generate them from authoritative models where practical. \\
Project models: typed registries, JSON Schema, OpenAPI~\citep{jsonSchemaDocumentation,openapi2026spec} & Schemas define data constraints; OpenAPI describes HTTP interfaces. Project declarations need explicit meanings for their relationships. & Keep models in ordinary project files; validate declarations and build the queries or derived views the work needs. \\
Clafer~\citep{claferDocumentation} & Declarative structural and feature models with constraints. & Generate instances, check consistency, and complete configurations. \\
BESSER~\citep{besserDocumentation} & A modeling framework with supported model families and transformations. & Create models visually, through textual imports, or in Python; generate code for supported targets. \\
SysML v2 / KerML~\citep{omg2025sysml,omg2025kerml} & Defined semantics for system modeling and its kernel constructs. & Connect interacting system views; select tools for the required editing, visualization, and analysis. \\
TLA+ and Alloy~\citep{lamport2002,tlaToolboxDocumentation,alloy6Documentation} & Formal behavioral specifications or relational and temporal models. & Use TLC or Alloy analysis under stated scopes and assumptions; TLAPS separately supports proofs. \\
\bottomrule
\end{longtable}
}

Choose the least elaborate approach that supports the questions and analyses you need. Agree on model-element meanings, authoritative sources, and change conventions. A rendering tool such as Mermaid can display a dependency graph, but the project must still define what its edges mean; generating that view from a maintained model can avoid a second account to reconcile. Agents can help construct models, generate views, and migrate representations when those conventions and checks are stable.

Some modeling approaches support formal verification: establishing mathematically that a model or implementation satisfies specified properties, subject to the assumptions and scope of the analysis. This can provide strong assurance, but only about the properties actually proved. Memory safety or the absence of certain illegal states may be comparatively straightforward to specify. Behavioral correctness is harder: before proving that a system behaves as required, engineers must determine what behavior is required. Agents can help formulate specifications and construct proofs, but the responsible engineer must understand and validate the properties, assumptions, and correspondence between model and implementation. A proof of a property that nobody has established as the right property is not evidence that the system meets its requirements. \autoref{3.2-models-and-the-semantic-gap:heading:agreement-is-not-correctness} develops this distinction between checking a property and establishing its adequacy.

The choice of modeling approach is also an economic decision. More explicit semantics can reduce the effort agents spend interpreting models, recovering missing relationships, and checking implementations against engineering obligations. In the terms of the cost-per-correct-admitted-change equation in \autoref{6.1-toward-a-theory-of-mage:heading:cost-and-failure-separately}, these savings must be weighed against the cost of constructing, validating, and maintaining the models and their tooling. Greater precision is worthwhile when its benefits across subsequent engineering work exceed those costs. It is not an objective in itself.

\subsection{Model the Change}\label{heading:use-a-change-as-the-first-model}

A structured engineering ticket can model the change you intend to make. It can preserve the intended behavior, rationale, acceptance criteria, affected concepts, constraints, related changes, and prior decisions. When it carries enough information to reason about one requested change, it functions as a change-scoped model. It should refer to longer-lived system knowledge rather than reproduce that knowledge in every ticket. For a defect, record the observed discrepancy and the intended correction; diagnosis may still be the difficult part of the work.

Identify who will use each link and for what decision. An agent may follow the ticket's component reference to the relevant ownership model; a reviewer may follow an invariant to its tests; a completion audit may compare the requested behavior with the landed result. Add a field when someone or something will use it. Fields without a consumer become additional material to keep current. The first useful investment may be a reliable link between an existing ticket, an architectural decision, and its acceptance evidence.

The ticket's lifetime matters. A ticket written in 2022 need not describe the product in 2026, but it can remain an accurate representation of the intent governing the 2022 change. Its persistence is useful history rather than a synchronization burden. A current architectural model makes a different claim and must be reconciled when the relevant system or intent changes. Mark that distinction before installing freshness checks: otherwise a well-intended audit can classify valid history as drift.

\subsection{Establish the Delegation Boundary}\label{heading:establish-the-delegation-boundary}

For the first trial, state what the agent may change, what evidence will support acceptance, and what should stop the work. Separate the producing agent's completion report from the evidence used to evaluate it. Existing tests, analyses, measurements, and review procedures may already provide much of that evidence. When they do not settle a consequential question, name the remaining judgment and who must exercise it. Keep an escalation path for a mistaken premise, an unavailable dependency, or an obligation the environment cannot adequately evaluate.

Run a small number of comparable work units and retain the interventions they required. Distinguish time spent creating new intent from time spent reconstructing known facts, checking routine conditions, or recovering from a recurring failure. These observations tell you what the next investment should address. A repeated escalation can indicate an appropriate human judgment boundary; it does not automatically justify a new evaluator or permission.

\subsection{From Local Practice to Shared Engineering Infrastructure}\label{heading:from-local-practice-to-shared-engineering-infrastructure}

MAGE can be applied at different scales. An engineer might begin with a model of one component, a behavioral specification for one feature, or a gate protecting one consequential obligation. Such investments are relatively inexpensive and can provide immediate benefits. Broader adoption requires more coordination: shared modeling conventions, stable component identities, common interfaces, and agreement about which obligations the engineering environment should enforce. The additional overhead can be substantial, but so can the return. When engineers and agents work from consistent models and controls, they can modify different components without repeatedly negotiating their relationships or rediscovering the rules governing their integration.

This tradeoff is familiar from conventional software engineering. Organizations standardize programming languages, test frameworks, build systems, and repository conventions partly to reduce the cost of collaboration and integration. A monorepo with uniform component ownership rules, for example, imposes conventions on individual teams but makes dependencies, ownership, and cross-component changes easier to manage. MAGE extends the same reasoning to engineering models and governance mechanisms. A shared architectural model can establish component boundaries, while common checks preserve them as different engineers and agents work independently. The benefit is not simply that each change becomes easier; it is that independently produced changes remain consistent with decisions made elsewhere in the system.

The appropriate scale depends on the expected return. A small project may need only a few local models and checks. A large organization may justify shared modeling infrastructure because its costs are distributed across many components, teams, and engineering episodes. Standardization can also impose unnecessary restrictions or make changes more expensive when conventions no longer fit. Expand the governed environment when the benefits of consistency, reuse, and coordination justify the additional construction and maintenance costs.

\subsection{Decide Whether to Expand}\label{heading:decide-whether-to-expand}

Expand the scope or consequence of delegated work when the evidence supports doing so. Check that someone remains able to inspect the result, diagnose failures, and continue essential work when the preferred agent or service is unavailable. Review the boundary by engineering surface; a successful workflow does not establish that the rest of the product is equally governable. This keeps the adoption decision connected to demonstrated capability and its costs.

Expansion is not the default. Once the selected capability has adequate representations, evidence, and enforcement, further machinery needs a new reason to exist. A useful first result may be modest: one dependency that stops being rediscovered, one recurring review decision made explicit, or one acceptance condition that no longer depends on memory.
